\documentclass[11pt,a4paper]{article}
\usepackage[utf8]{inputenc}
\usepackage[T1]{fontenc}
\usepackage[spanish]{babel}
\usepackage[a4paper,margin=1.8cm]{geometry}
\usepackage{xcolor}
\usepackage{tikz}
\usetikzlibrary{positioning,arrows.meta,shapes.geometric,fit,backgrounds}

\definecolor{rel}{RGB}{33,102,172}
\definecolor{ts}{RGB}{33,150,120}
\definecolor{ml}{RGB}{190,100,40}
\definecolor{app}{RGB}{80,80,80}

\tikzset{
  box/.style={draw, rounded corners=2pt, align=center, minimum height=0.8cm,
              minimum width=3.1cm, font=\small, inner sep=4pt},
  lbl/.style={font=\scriptsize\itshape, text=gray}
}

\title{\textbf{Arquitectura propuesta}\\[0.3em]
\Large Predicción de demanda espaciotemporal (TimescaleDB + PostGIS)}
\author{García, Barreto, Chulluncuy, Hilario}
\date{Base de Datos (MCD8008) -- Parte 2}

\begin{document}
\maketitle

\begin{center}
\begin{tikzpicture}[
  node distance=0.55cm and 1.1cm,
  arr/.style={-{Stealth[length=2.2mm]}, thick}
]

% Aplicaciones y API
\node[box, fill=app!10] (app) {APP Web / Mobile};
\node[box, fill=app!10, below=of app] (api) {API};

% Rama relacional
\node[box, fill=rel!15, below left=of api, xshift=-1.6cm] (pg) {PostgreSQL\\ \scriptsize (core transaccional)};
\node[box, fill=rel!15, below=of pg] (datos) {Datos de negocio};

% Rama analitica
\node[box, fill=ts!15, below right=of api, xshift=1.6cm] (ts) {TimescaleDB + PostGIS\\ \scriptsize (series + espacial)};
\node[box, fill=ts!15, below=of ts] (fs) {Feature Store};
\node[box, fill=ml!15, below=of fs] (ml) {ML Pipeline\\ \scriptsize (entrenamiento)};
\node[box, fill=ml!15, below=of ml] (reg) {Model Registry\\ \scriptsize (MLflow)};
\node[box, fill=ml!15, below=of reg] (inf) {API de inferencia\\ \scriptsize (FastAPI)};
\node[box, fill=ml!15, below=of inf] (out) {Dashboard +\\ alertas de anomalías};

\draw[arr] (app) -- (api);
\draw[arr] (api.south) |- (pg);
\draw[arr] (api.south) |- (ts);
\draw[arr] (pg) -- (datos);
\draw[arr] (ts) -- (fs);
\draw[arr] (fs) -- (ml);
\draw[arr] (ml) -- (reg);
\draw[arr] (reg) -- (inf);
\draw[arr] (inf) -- (out);

\end{tikzpicture}
\end{center}

\section{Descripción de componentes}

\begin{itemize}
  \item \textbf{APP Web / Mobile y API:} aplicaciones consumidoras del pronóstico y de las
        alertas.
  \item \textbf{PostgreSQL (core transaccional):} mantiene viajes, catálogos, lotes de ingesta,
        reglas de calidad e integridad referencial. Es la fuente de verdad.
  \item \textbf{TimescaleDB + PostGIS:} capa analítica con el hypertable
        \texttt{demanda\_serie}, el continuous aggregate diario y la columna espacial
        \texttt{geo} con índice GiST.
  \item \textbf{Feature Store:} materializa las variables (rezagos, medias móviles, variables
        espaciales) para entrenamiento y para inferencia en línea.
  \item \textbf{ML Pipeline:} entrena el modelo (baseline LightGBM) con validación temporal
        \emph{walk-forward}.
  \item \textbf{Model Registry (MLflow):} versiona y sirve los modelos entrenados.
  \item \textbf{API de inferencia (FastAPI):} expone el pronóstico horario por zona y
        plataforma con latencia objetivo $<100$\,ms.
  \item \textbf{Dashboard + alertas:} visualización de demanda y detección de anomalías.
\end{itemize}

\section{Flujo de datos}

\begin{enumerate}
  \item La ingesta mensual carga los viajes en PostgreSQL (núcleo transaccional).
  \item El agregado \texttt{demanda\_hora} se proyecta al hypertable
        \texttt{demanda\_serie} (TimescaleDB) y se enriquece con la geometría de
        \texttt{zona\_taxi} (PostGIS).
  \item El feature store consolida variables temporales y espaciales.
  \item El pipeline entrena el modelo y lo registra en MLflow.
  \item La API de inferencia sirve pronósticos y el dashboard consume el resultado,
        disparando alertas ante anomalías.
\end{enumerate}

\end{document}
